\chapter{Considerações Finais}


A etapa de simulação do MUSE foi validada no robô Go1 no simulador Gazebo, e atualmente estão sendo feitos testes no robô Go1 real. Espera-se que, com o retorno para o Brasil, seja possível trabalhar mais diretamente no robô Go2.

A metodologia escolhida, que combina simulação com o monitoramento de contato e de deslizamento, tem se mostrado essencial para o sucesso do projeto. O uso do Gazebo permitiu ajustar a arquitetura de observadores em cascata (NLO e XKF) em um ambiente seguro, minimizando os riscos antes de levar o sistema para o hardware real. Esse preparo é fundamental para garantir que o estimador lide bem com o escorregamento em superfícies de baixo atrito, permitindo ajustes finos que seriam muito mais difíceis de realizar diretamente no robô.

Por fim, o trabalho abre portas para diversas aplicações futuras, oferecendo uma base sólida para o grupo MOBILab da UDESC. Os próximos passos podem focar no aprimoramento da fusão sensorial e na avaliação do sistema sob condições ainda mais dinâmicas.
% \section{Conclusões}


% \section{Limitações}


% \section{Trabalhos Futuros}